% Expository adaptation, October 7, 2026. Original mathematics: OpenAI, September 23, 2026.
% Prose revised and explanations expanded; see README.md for the source and scope.
\section{Comparing vertices by short fingerprints}
\label{sec:fp-blocks}

Two descriptions may reach the same destination. Their contributions must be combined, but comparing full identifiers can require too much information to remain stored during recursion. We therefore use short hash values, or fingerprints, when the budget decrease is small. The vertex universe has size $2^{O(B)}$, whereas the fingerprint has only $B$ possible values.

A collision can change the subsequent computation, including the destinations queried later. It would be incorrect to apply a fixed-pair collision bound directly to these adaptively chosen pairs. The proof instead compares the actual computation with one using exact equality throughout one block of budgets. This comparison supplies a set of candidate vertices that does not depend on the current block's hash seed. We can bound collisions on that set and charge departures from it to errors in lower calls.

There is a second reason to define the comparison carefully. Its conditional averages must be deterministic at the queried pair, as required by the compression proof. Freezing a random coordinate on which the sample still depended would invalidate that requirement. The rules below specify both which hash seeds a task may inspect and which ones its averaging walk varies.

We first construct the fingerprints and explain how to compare them without retaining an index at every call. We then give the exact update rule and prove the adaptive error bound.

We use the root error metric $\mathfrak d_{u,a}$ for bounded nonnegative row
outputs, its ordinary version $\|\cdot\|_{u,a;s}$ for raw arrays, and the
ordinary vertex metric for vectors.  For a row quantity $Z(x,\sigma)$ define
also
\begin{equation}
 \|Z\|_{\mathrm{row},u;s}^s
 =\frac1N\mathbb E_\sigma\sum_x\frac{X_u(x)}u|Z(x,\sigma)|^s.
 \label{eq:fp-row-measure}
\end{equation}
The value of $s$ remains $8$.  Supremums over endpoints below are over actual
ID positions, with the appropriate retained-domain masks; virtual default
positions are never counted with the multiplicity of the ambient universe.

\subsection{Choose one short hash seed for each budget block}

The fixed-pair estimate below combines polynomial evaluation, as in
classical algebraic fingerprinting~\cite{Schwartz1980}, with affine
universal hashing~\cite{CarterWegman1979}. Its proof is included to
specify the fair-bit distribution and space bound. The adaptive use of
these fingerprints requires the separate block comparison that follows.

\begin{lemma}[Short fingerprints]\label{lem:fp-hash}
Suppose every endpoint has an injective, fixed-length binary encoding of
length $L\le c_{\mathrm{ID}}B$.  There is a uniform construction from
$O(\log B)$ independent fair bits of a map
$k:\{\text{endpoint IDs}\}\longrightarrow\{1,\ldots,B\}$ such that, for
distinct fixed IDs $v,w$,
\begin{equation}
 \Pr[k(v)=k(w)]\le C/B\le C B^{-1/2}.
 \label{eq:fp-pair-collision}
\end{equation}
Evaluating the map at a presented endpoint needs $O(\log B)$ arithmetic
scratch in addition to access to its encoding.
\end{lemma}

\begin{proof}
Pad all encodings to the same length and put
$P_v(T)=\sum_{i=0}^{L-1}v_iT^i$.  Choose a prime
$B^3\le q\le2B^3$, increasing the fixed lower bound on $B$ if necessary.
Thus $L<q$.  Let $\ell=\lceil\log_2q\rceil$.  Split a string of $3\ell$
fair bits into independent integers $R,A,C$ in $\{0,\ldots,2^\ell-1\}$,
and reduce them modulo $q$ to obtain $r,\alpha,\beta$.  Each residue has
probability at most $2/q$.  Let
$\tau:\mathbb F_q\longrightarrow\{1,\ldots,B\}$ partition the standard
residues into consecutive sets of sizes differing by at most one, and set
\[
 k(v)=\tau\bigl(\alpha P_v(r)+\beta\bigr).
\]
All computations are in $\mathbb F_q$ before applying $\tau$.

The nonzero polynomial $P_v-P_w$ has degree at most $L-1$, so it has at most
$L-1$ roots.  This root bound follows by repeatedly factoring $T-r_0$ at a
root; it applies here because the coefficient difference $1$ or $-1$ is
nonzero in $\mathbb F_q$.  Consequently
\[
 \Pr[P_v(r)=P_w(r)]\le 2(L-1)/q.
\]
For distinct evaluations $z,z'$, the map
$(\alpha,\beta)\mapsto(\alpha z+\beta,\alpha z'+\beta)$ is a bijection
of $\mathbb F_q^2$.  Under uniform parameters its two outputs are uniform
and independent.  The probability that their $\tau$-values agree is
\[
 q^{-2}\sum_{j=1}^B|\tau^{-1}(j)|^2
 \le \frac{\lceil q/B\rceil}{q}\le\frac1B+\frac1q.
\]
The actual joint law of $(\alpha,\beta)$ is pointwise at most four times
the uniform law and is independent of $r$.  Adding the two cases gives
$2(L-1)/q+4/B+4/q=O(1/B)$.  A polynomial evaluation by Horner's rule and
the residue arithmetic use $O(\log q)=O(\log B)$ bits.  The prime can be
chosen uniformly by searching and trial division in that space.  The
existence of a prime in the indicated interval is the same elementary
prime-interval fact used in choosing the environment field.
\end{proof}

A new independent hash seed for every recursive call would consume too much randomness. Instead, divide the statistical budgets into blocks and share one seed, called a salt, within each block. A block has width proportional to $\log B$, and a salt has $O(\log B)$ bits, so the total over $O(B/\log B)$ blocks is $O(B)$ bits.

Fix a small positive rational $\epsilon$, to be chosen last.
This coefficient controls the fingerprint block width. Write
\begin{equation}
 h_*=\lfloor\epsilon\log_2 B\rfloor,\qquad
 M_b=bh_*,\qquad
 \mathcal I_b=\{M_b,\ldots,M_b+h_*-1\}.
 \label{eq:fp-blocks}
\end{equation}
Enlarge the fixed minimum $B$ so $h_*\ge1$ and
$h_*\ge(\epsilon/2)\log_2B$.  Give each block meeting
$\{0,\ldots,M_{\mathrm{top}}\}$ an independent salt from
Lemma~\ref{lem:fp-hash}.  Pad each salt to a multiple of the fixed integer
$H$ by unused independent bits.  Since $M_{\mathrm{top}}=O(B)$, the total
number of bits is
\begin{equation}
 O\bigl((1+B/h_*)(\log B+H)\bigr)=O(B).
 \label{eq:fp-salt-space}
\end{equation}
Here and throughout, constants may depend on the fixed $\epsilon$ and $H$.
We enumerate the original fair-bit strings when enumerating environments;
reducing chunks modulo $q$ does not change that sampling convention.

The restriction on which salts a task reads is essential to the comparison argument: lower-budget blocks must not inspect the current block's salt. The following access rule is part of the algorithm.  A task whose budget
lies in $\mathcal I_b$ may inspect salts with indices at most $b$.  Its
conditional walks vary all earlier salts as free bits and freeze the salt
of block $b$ and every later salt.  The restriction applies to structural
choices, port orders, validity tests, and predecessor tests as well as to
numerical values.  The controller construction below implements this rule
for every access operation, including inverse queries.  It is not enough
to impose it only on the final output values.

This rule is stated in the coordinates of the task's \emph{effective
environment}.  Write $\sigma$ for that environment when analyzing a
task.  A fixed ancestor variant may mix its free salt coordinates before
supplying $\sigma$ to a descendant; the descendant's table depends on
the resulting environment, not on the word that produced it.  Thus the
rule does not forbid reading a later-indexed bit of the original master
environment while replaying that word.  It requires the descendant's
values, port orders, and access operations to be functions only of
its permitted effective coordinates.  Lemma~\ref{lem:ctrl-replay}
implements this distinction.  Each fixed variant is measure preserving
by Lemma~\ref{lem:env-variants}, so the integrated estimates may change
variables to its effective environment.  We do not assume that an
adaptively selected variant leaves a uniform conditional law.  At a
block-$b$ task, its own variants act only on effective salts earlier
than $b$ and freeze its effective block-$b$ salt and all later ones.

\subsection{Compare keys while restoring a shared bit vector}

A short key is not yet enough: keeping a $\log B$-bit index at each of many recursive levels could still exceed the space bound. Instead, use one shared $B$-bit vector as temporary workspace. A comparison may alter this vector during its execution, but it must restore exactly the contents it received. This is called catalytic use of memory; no free or uncounted memory is being assumed.

The idea is to flip the bit indexed by the first key and read the bit indexed by the second key both before and after undoing that flip. The two reads differ exactly when the keys agree. Each key is recomputed only when its vertex is visited, so no key index must survive a recursive excursion.
The cancellation of arbitrary initial contents also appears in
transparent programs for catalytic computation
\cite[Section~3]{BuhrmanCleveKouckyLoffSpeelman2014}.  Here the shared bit
vector is included in the ordinary work-space bound.

\begin{lemma}[Restoring equality test]\label{lem:fp-catalytic}
Let $V\in\{0,1\}^B$ be a shared vector with arbitrary initial contents.
Suppose an excursion to either of two deterministically specified
endpoints returns the cursor, and its path and result are independent
of the incoming contents of $V$.  Such excursions may use nested
procedures which restore $V$.  There is a procedure which xors the bit
$\mathbf1_{k(v)=k(w)}$ into a local result bit, restores $V$, and needs
only constant local data across four endpoint excursions.
\end{lemma}

\begin{proof}
Starting with result bit $r_0$, perform four restoring excursions in this order:
\begin{enumerate}
\item Visit $v$ and flip $V_{k(v)}$.
\item Visit $w$ and xor $V_{k(w)}$ into the result.
\item Visit $v$ again and flip $V_{k(v)}$ back.
\item Visit $w$ again and xor $V_{k(w)}$ into the result.
\end{enumerate}
The key is evaluated only at the vertex currently visited. Across the excursions, retain only the phase and result bit.  If $V^0$ is the incoming vector, the first
read equals $V^0_{k(w)}\mathbin\oplus\mathbf1_{k(v)=k(w)}$ and the
second equals $V^0_{k(w)}$.  Thus the final result is
$r_0\mathbin\oplus\mathbf1_{k(v)=k(w)}$ and both flips restore $V^0$.

For nesting, use induction on the well-founded procedure order.  A
nested call sees some arbitrary incoming vector, restores precisely
that vector, and returns a result independent of it.  Hence the four
outer endpoint paths and their keys are the same on repeated visits,
even between the two outer flips.  The terminal evaluation, read, or
flip is a local operation with no further structural call.  Only a
phase and a result bit must survive across the excursions; a key index
is recomputed at each visit in shared scratch.  The result bit being
accumulated is not used to choose either endpoint path.  These facts
give the induction invariant and the claimed local storage bound.
\end{proof}

For $d>2h_*$, exact bit-by-bit ID comparison is retained.  A bit position
then takes $O(\log B)=O(d)$ suspended bits, with constants allowed to
depend on the already fixed $\epsilon$.  For small slack,
Lemma~\ref{lem:fp-catalytic} avoids that per-frame index.  The shared
vector itself occupies $B$ bits and, by the lemma, may be reused by all
nested comparisons.  This algebra establishes restoration; the
controller and traversal section supplies the deterministic excursion
and call-order interfaces needed to apply it.

The endpoint action evaluates the key with the invoking comparison's
effective salt, even when the excursion visits a lower-budget table.
It makes no further access query and does not choose or alter the
excursion.  Thus the lower path remains independent of later salts.
Lemmas~\ref{lem:ctrl-equality} and~\ref{lem:ctrl-replay} implement this
separation for all restoring accesses, including inverse queries.


\subsection{The actual update and the exact-equality comparison}

We now specify the destination-grouping step in the estimator description. The construction is the same row update already outlined in Section~\ref{sec:estimator-description}; here the treatment of collisions is made exact. Inactive tasks still return zero outputs and empty access classes.  For an active
row task, the additive schedule of Section~\ref{sec:est} supplies
candidate increments from smaller-budget outputs.  Process the groups in
decreasing slack.  In each group, distribute increments by the rule below,
then use the compression, range clamps, and shortlist rule of
Section~\ref{sec:buf-compression}.  After the last group, form the
pre-incoming row functions and precursor classes from the completed arrays,
then assemble any incoming gates and column multipliers as in
Section~\ref{sec:atom}.  The separate column-vector calls have smaller
budget.  Vector tasks use the same schedule to sum their scalar terms,
with their prescribed range clamps and final clipping; they require no
endpoint grouping or row-list compression.

The important choice is to group only the new increments. Combining old stored values when their keys collide could introduce an error as large as those values, even if the new increments were tiny. Keeping the old values at their existing vertices allows the collision cost to be measured using the small successive-budget differences.

The row-group distribution rule is as follows.
For a dyadic group with slack $d>2h_*$, compare endpoint IDs exactly.  For a
group with $d\le2h_*$, use the key from the task's block as follows.
First discard invalid candidates, including invalid guessed wraps,
individually.  Sum all new candidate increments having the same key.
Add a bucket's sum to every old representative with that key.  If the
bucket has no old representative, introduce exactly its first valid new
candidate, with its previous value equal to the old default, and add the
bucket sum there.  Do not combine or discard old representatives in this
key step.  Apply the same full-array compression, clamping, and shortlist
rule as in the preceding section.

Distinct actual representatives remain distinct.  Indeed two introduced
representatives have different keys, and an introduced representative's
key differs from every old key.  Increments at one and the same actual
endpoint are always combined, even when other endpoints collide with it.
Stored values are never summed by key.  The two channels use the same keys
and the same representative selection for the entire group.

To analyze this adaptive rule, remove the current block's dependence on its own salt. For an analytical comparison at a budget $M\in\mathcal I_b$, define the
\emph{clean calculation for block $b$} by replacing key comparisons by
exact ID comparisons throughout that block, recursively in every child
whose budget still belongs to $\mathcal I_b$.  Children below $M_b$ are
the actual algorithms, with exactly the same environments in the two
calculations.  This replacement applies to the entire numerical and
structural calculation: candidate validity, representative selection,
compression, exposed classes, and their port orders.  Descriptor
enumeration uses fixed index orders.  Each finite-accuracy test for a
shortlist or exposed class retains its prescribed accuracy level and
tie convention, independently of any requested output accuracy.  All
these operations use the corresponding clean operands.  Thus the clean
calculation specifies its access choices as well as its represented
arrays.

Denote the actual and clean versions by superscripts $\mathrm{ac}$ and
$\mathrm{cl}$.  Both still perform compression and clamping.  The clean
calculation is a finite analytical recursion; its exact comparisons are
not asserted to satisfy the eventual space bound.

Clamping is part of both definitions.  The clean error reserve uses
Lemma~\ref{lem:buf-clamping} for exact-ID increments, which are
independent of the stored exception list.  That lemma preserves the
relevant row functions and witnesses; it need not preserve the list or
port presentation of an unclamped calculation.  The coupling below
compares the two clamped calculations directly.  Its propagation bound
uses the Lipschitz compression map and the nonexpansiveness of each
common clamp, without introducing an unclamped actual calculation.

\begin{lemma}[Independence of the clean block]\label{lem:fp-clean}
Every clean value and structural choice for block $b$ is independent of
that block's effective salt and of later effective salts.  In a task of
budget $M\in\mathcal I_b$, the actual and clean row states immediately before
the first group with $d\le2h_*$ are identical and independent of that
salt.  Their stored lists at this boundary have size $2^{O(h_*)}$.
The deterministic slice-mean assertion required by
Lemma~\ref{lem:est-prefix} holds for the clean calculation.
\end{lemma}

\begin{proof}
We must prove independence for the entire computation, including which ports it exposes; independence of the final numerical answer alone would not be enough. Induct on the budget and, within a task, on its prescribed sequence of operations.  Lower-block algorithms ignore the salt of block $b$ and
later salts by the access rule; same-block clean children ignore them
by induction.  Candidate tests use these children, exact endpoint
grouping precedes compression, and the compressed row determines its
exposed classes and pre-incoming scores.  Incoming ranks use those
completed scores.  Separate column-vector queries likewise use a
smaller-budget clean child in this block or an actual child below it.
All statements here concern the current effective environment: each
same-block transformation preserves its current and later salts, and
each lower-block transformation acts on a subset of its earlier salts.
Each numerical recipe uses the same fixed operand
order, and each structural test uses its fixed accuracy level and tie
rule.  Its result therefore depends only on earlier clean operands,
including their finite numerical approximations.  The fixed index
orders then give salt-independent representatives and ports.

In particular, exact full-array addition, compression, and the prescribed
row functions are salt independent.  Reordering paths or choosing a
different representative of the same endpoint would not change an exact
sum; retaining an extra default-valued position would not change the
represented array.  The fixed conventions above also remove these
presentation choices from the structural calculation.

For a group with $d>2h_*$, every child budget is at most $M-d<M_b$.
Both runs use the same lower-block children and exact comparisons there.
All their list operations, including their deterministic tie conventions,
ignore the present salt.  Thus their boundary states agree.  The last
larger dyadic slack, if there is one, is at most $4h_*$, so its shortlist
has at most $2^{O(h_*)}$ positions.  If there is no larger group, the
boundary list is empty.

Finally, the clean estimand depends only on rank channels and earlier
salts.  These are precisely the variables averaged by its conditional
walk, subject to its stated vertex conditions.  Same-block salts are
irrelevant, and lower-block salts are free variables of this walk.
Projection onto the uniform measure of the slice therefore gives a
deterministic quantity at its indicated pair or vertex, with the
prescribed retained-domain masks.  In particular, no random dependence
on a frozen salt remains in this mean.  For the second channel, its
deterministic mean magnitude is extended by its prescribed narrower
endpoint mask, as required in Section~\ref{sec:buf-compression}.
The salt rule is required of all access operations, including inverse
queries and port orders; Lemma~\ref{lem:ctrl-replay} supplies its
implementation.  Lemmas~\ref{lem:env-gap} and \ref{lem:env-variants}
therefore apply to this function.  This verifies the deterministic-mean
hypothesis of Proposition~\ref{prop:est-clean}; it does not assert
independence between errors of different children.
\end{proof}

\subsection{Find a candidate set independent of the current hash seed}

The exact-equality comparison determines a candidate set independent of the current salt. We enlarge it slightly using the smooth support witnesses, so that if an actual access leaves the set, some child witness must have changed from zero to one. This turns a discrepancy in access choices, even at a zero-weight entry, into a numerical error already covered by the induction.

We recall the relevant interfaces from Lemma~\ref{lem:atom-ports}.
Every exposed actual path port has a broader smooth witness equal to one
at that port.  Each witness has bounded row support.  The actual
correction ports, including exposed ports of weight zero, have incoming
multiplicity at most $K_c$.  Tables contain one entry per pair of actual
IDs.  After fixing a copy part, the child tuple is common over the whole
table; in particular, a right child reached from a midpoint does not
depend on the discarded old start.  All these statements hold for
arbitrary computed raw arrays, not just for accurate ones.

The incoming cap applies to the entire exposed correction relation,
including its zero-weight entries.  The broader auxiliary score class
needs only a row bound: it is queried for scores and witnesses, whereas
the left factor of a two-factor path is a capped correction table.
This distinction permits the source sum below to be removed even when
an offending candidate has zero numerical weight.

Let $\mathcal G_b$ be the sigma-field generated by all coordinates of the
task's effective environment except the salt of block $b$.  Fix a row
$x$ and condition on $\mathcal G_b$ with $X_u(x)=1$.  Take its boundary list from
Lemma~\ref{lem:fp-clean}, including every stored exact-default slot.
For every small group and every immediate
descriptor variant, evaluate the clean factors and their witnesses in
the environment obtained by that fixed variant.  This transformation
leaves the present and later salts frozen.  Include the endpoints in
each direct witness support and, for a two-factor path, the endpoints
in the right witness supports at all midpoints in the left witness
support.  Include local base endpoints when applicable.  Let
$\mathcal E_x$ be the union of these endpoints and the boundary list.
Set $\mathcal E_x=\varnothing$ when $X_u(x)=0$.
Unused or invalid clean descriptors may enlarge this set.  Actual
candidates must still satisfy their wrap tests and all transformed
source, midpoint, and terminal memberships; every error estimate below
retains the memberships needed for its child metrics.

\begin{lemma}[Size and independence of the envelope]
\label{lem:fp-envelope}
The set $\mathcal E_x$ is independent of the present salt and has size at
most $2^{C_Eh_*}$ for a fixed $C_E$.  All clean nondefault representatives
in the small groups belong to this set.  Consequently, conditional on
the other data,
\begin{equation}
 \Pr[\text{two distinct IDs in $\mathcal E_x$ have the same key}]
 \le C\,2^{2C_Eh_*}B^{-1/2}.
 \label{eq:fp-envelope-collision}
\end{equation}
\end{lemma}

\begin{proof}
First fix every coordinate except the current salt. This also fixes the comparison computation and all of its witness supports. Each fixed variant is independent of the present salt and leaves it
unchanged.  Lemma~\ref{lem:fp-clean} therefore makes its clean witness
supports $\mathcal G_b$-measurable.  We use these supports as sets of
actual IDs, not the order of a port enumeration.  At a group of slack
$d$ there are $2^{O(d)}$ immediate descriptors, including variants,
difference parts, and copy choices.  The number of endpoints contributed by a direct
witness is bounded.  A two-factor path contributes a bounded number of
left witness midpoints and, for each, a bounded number of right witness
endpoints.  Summing over the dyadic groups with $d\le2h_*$ and adding the
boundary list proves the size bound.  No incoming bound on a clean
witness is used here.

Clean exposed ports lie in their clean witness supports.  Apart from the
boundary representatives, exact updates can therefore create nondefault
values only at the included endpoints.  Compression maps all untouched
default positions to the new default.  Any extra exact-default slots in
a shortlist may be omitted when describing its full array.  This proves
the representative claim.  Equation~\eqref{eq:fp-envelope-collision}
is now the pair union bound from Lemma~\ref{lem:fp-hash}, conditional on
the fixed envelope.
\end{proof}

Let $\mathcal F_x$ be the event that some valid actual candidate in the
small groups has endpoint outside $\mathcal E_x$, whether or not its
increment is zero.  As representatives are
only carried forward or chosen among new candidates, all actual
representatives stay in the envelope outside this event.
This includes default-valued representatives: they start in the full
boundary list or inherit a candidate endpoint.  Dummy copies used to
compute ranks are not new real representatives.  Thus omitting an
exact-default clean slot when describing its full array does not omit
an actual address from the collision analysis.

\begin{lemma}[Charging an escaped candidate]\label{lem:fp-escape}
Suppose actual--clean child discrepancies at budget $m$ have their
prescribed metric bounds $2^{-p_m}\lambda_m$.  The contribution of
$\mathcal F_x$ to any bounded final pre-incoming row output has norm at
most
\begin{equation}
 2^{-p}\sum_{\substack{d\le2h_*\\m\in\mathcal I_b,\ m\le M-d}}
                 2^{C d}\lambda_m,
 \label{eq:fp-escape-charge}
\end{equation}
where the sum includes the finitely many dependency types and their
descriptor multiplicities in its constant $C$.  Equivalently, one may
use the maximum of $\lambda_m$ over the children of each slack group
and increase $C$.
\end{lemma}

\begin{proof}
Locate the first factor of an actual path that leaves the corresponding clean witness support. At that factor, the actual witness is one and the clean witness is zero. Thus an access disagreement produces a unit discrepancy, whether or not the path's numerical contribution vanishes.

An actual endpoint outside the envelope has an offending path factor.
For a direct path its actual port is outside the clean witness support.
For a two-factor path, either its actual left port is outside its clean
witness support, or its midpoint lies in that support and its actual
right port is outside the corresponding clean right witness support.
At the offending port the actual witness is one and the clean witness
is zero.  Thus the indicator of this event is at most the $s$-th power
of their root discrepancy there.  Union bounding over candidate paths
gives a sum of such nonnegative powers.  The unit-witness property holds
also for exposed zero-weight ports; no lower bound on the candidate's
increment is being assumed.

Once such a path occurs, we charge the whole bounded row-output error to its occurrence. Charging only the path's increment would be insufficient: even a zero increment can affect the selection of representatives. Each output under consideration is bounded and has bounded row support.
The union of the two output supports is still bounded, by
Equation~\eqref{eq:buf-row-support}.  Hence its entire row error power on
$\mathcal F_x$ is at most a fixed multiple of
$\mathbf1_{\mathcal F_x}$.  We charge this row-failure indicator, rather
than the size of the offending increment.  There is no endpoint sum left
on the output side to pay for.

The right-factor charge is the case where the old start must be summed
out.  Fix a guessed variant and a two-factor path $x\to z\to y$.
Change variables by its measure-preserving transformation, as in
Lemma~\ref{lem:env-variants}; in the following display all memberships,
ports, and witness values are evaluated in that transformed environment.
Write $e_R(z,y)$ for the root discrepancy of the two right witnesses and
$\mathcal P^{\mathrm{ac}}_L$ for the actual left exposed-port relation.
The candidate's validity retains $X_u(x)$, $Z_f(z)$, and the right
child's endpoint membership $Y_{a'}(y)$.  Dropping only surplus guards,
the right-port contribution is at most a constant times
\begin{align}
 &\frac1N\mathbb E\sum_x\frac{X_u(x)}u
       \sum_{z:(x,z)\in\mathcal P^{\mathrm{ac}}_L}Z_f(z)
       \sum_{y:Y_{a'}(y)=1}|e_R(z,y)|^s \notag\\
 &\qquad\le\frac{K_c}{N u}\mathbb E\sum_z Z_f(z)
       \sum_{y:Y_{a'}(y)=1}|e_R(z,y)|^s
   =K_c\frac f u\,\mathfrak d_{f,a'}(W_R^{\mathrm{ac}},
                                      W_R^{\mathrm{cl}})^s.
 \label{eq:fp-actual-incoming}
\end{align}
The middle inequality uses the incoming cap of the \emph{actual} left
port relation.  It does not use a clean incoming cap.  The right task's
tuple is independent of $x$, so the same discrepancy is being summed
when the old start is removed.  This condition is necessary for the
displayed equality with the right metric.
For the fixed variant and copy part, the right tuple has only the
prescribed task parameters and transformed environment.  The midpoint
$z$ is its vertex argument, and the old source $x$ is not a parameter.
Its final gate and column queries use that same effective tuple, rather
than the ancestral descriptor that happened to discover the entry.

For a direct or left offending port, sum the child's witness discrepancy
over its own row and endpoint domain; the bounded number of possible
right ports incurs only a fixed factor.  Projected paths are partitioned
by the fixed copy symbols, with constant multiplicity.  Normalization
and preparation dependencies make the same direct witness comparison.
If a candidate arose from a narrower terminal slice, its validity mask
may be dropped in this row-event estimate.  Thus no factor
$\pi_a/\pi_t$ is incurred.  In contrast, the transformed source hit
cannot be omitted before the change of variables and the incoming-cap
sum; it has been retained in Equation~\eqref{eq:fp-actual-incoming}.

At slack $d$, scale ratios, variants, and immediate path counts cost at
most $2^{O(d)}$, as in Lemma~\ref{lem:est-raw-difference}.  The possible
reset factor in Equation~\eqref{eq:fp-actual-incoming} is absorbed by the
precision gain.  Admissible row tasks have $u\le f$, as specified in
Section~\ref{sec:atom}.  Equation~\eqref{eq:est-child-precision} gives,
after the scale-dependent $O(d)$ terms,
\begin{align*}
 p_m&\ge p-O(d)+\tfrac12\log_2(f/u),\\
 (f/u)^{1/s}2^{-p_m}
 &\le 2^{-p+O(d)}(f/u)^{1/8-1/2}
 \le 2^{-p+O(d)}.
\end{align*}
Take $s$-th roots and use Minkowski to sum the charges.  A child below
$M_b$ is identical in the two versions and contributes zero.  All
remaining children satisfy $m\le M-d$ and belong to $\mathcal I_b$,
which proves Equation~\eqref{eq:fp-escape-charge}.
\end{proof}

\subsection{Bound the additional error from collisions}

The escape estimate leaves rows whose candidates stay in
$\mathcal E_x$.  On these rows we separate child discrepancies from
collisions among the clean endpoints.  We compare full arrays with
their defaults, so different shortlist presentations require no extra
error term.  The following observation explains the key update rule.

\begin{lemma}[Only increments are aliased]\label{lem:fp-buckets}
Consider one small group on $\mathcal F_x^c$.  Let
$\Delta_{\mathrm{ac}}(y),\Delta_{\mathrm{cl}}(y)$ be its unaliased
increments, namely its true-ID sums using actual and clean children,
respectively.  Let $k_d\le2^{C d}$ bound its number of immediate
descriptors, and let $\mathcal C_x$ denote a collision on $\mathcal E_x$.
The discrepancy between the actual bucket-distributed increment and
the clean true-ID increment is at most
\begin{equation}
 (k_d+1)\left(
       \sup_y|\Delta_{\mathrm{ac}}(y)-\Delta_{\mathrm{cl}}(y)|
       +\mathbf1_{\mathcal C_x}\sup_y|\Delta_{\mathrm{cl}}(y)|
                  \right)
 \label{eq:fp-bucket-bound}
\end{equation}
in sup norm, simultaneously for both channels.  The old full-array
discrepancy is added to this bound with coefficient one.
\end{lemma}

\begin{proof}
We compare the arrays before compression, separating a change in the children from a collision among the clean destinations. When all envelope keys are distinct, the rule gives exactly the actual
true-ID increment: each old occurrence is updated at its own ID, and
each genuinely new ID has its own bucket.  Thus the discrepancy is
simply $\Delta_{\mathrm{ac}}-\Delta_{\mathrm{cl}}$.

When there is a collision, each updated representative receives a sum
of increments at at most $k_d$ distinct actual addresses.  Its magnitude
is at most $k_d\sup_y|\Delta_{\mathrm{ac}}(y)|$.  At positions not
updated it is zero.  At either type of position subtracting the clean
increment costs at most a further
$\sup_y|\Delta_{\mathrm{cl}}(y)|$.  Apply
\[
 \sup_y|\Delta_{\mathrm{ac}}(y)|
 \le\sup_y|\Delta_{\mathrm{ac}}(y)-\Delta_{\mathrm{cl}}(y)|
      +\sup_y|\Delta_{\mathrm{cl}}(y)|
\]
to obtain Equation~\eqref{eq:fp-bucket-bound}.  This also covers an
increment duplicated at several old representatives and a missing
increment at an unchosen new representative.

At no point are old values combined or moved to a different ID.  An
unlisted old position has exactly the old default, and a newly selected
representative starts with that value.  Therefore the full-array
comparison before compression is the old discrepancy plus the increment
discrepancy just bounded.  There is no term involving the height of an
old value multiplied by a collision indicator.
\end{proof}

In applying this lemma to a group, combine contributions at equal actual
endpoints \emph{within each paired difference} before measuring its
clean height.  At a fixed endpoint and original generator word, the
outer slice and wrap tests are the same for budgets $m$ and $m-1$.
Their constituent port sets may differ, but each set contains all
nonzero entries and enumerates each actual ID pair once; an absent
entry contributes zero.  Their true-ID sums are therefore exactly the
two terms of the paired difference.  Subtracting the same stochastic
true-atom term from both sides, as in
Lemma~\ref{lem:est-raw-difference}, preserves its cancellation.
One can then sum the resulting sup bounds over difference layers.
This remains legitimate
if a first-bucket representative is chosen across several layers:
the proof of Lemma~\ref{lem:fp-buckets} applies to their total increment,
and the triangle inequality bounds the total by the sum of the paired
bounds.  A descriptor with zero increment contributes only to $k_d$ in
this bound.  It may affect representative selection, which is already
covered by Lemma~\ref{lem:fp-buckets} and the escape event.

The remaining probabilistic step must be applied to quantities independent of the current salt. The clean increment heights and the enlarged candidate set have this property. The event that the actual computation stays inside that set need not have it, so we do not condition on that event when applying the collision bound.

\begin{lemma}[Direct cost of collisions]\label{lem:fp-direct}
The direct collision contribution to any final pre-incoming output has
norm at most
\begin{equation}
             2^{-p}\,2^{C h_*}B^{-1/(2s)},
 \label{eq:fp-direct-charge}
\end{equation}
for a fixed $C$ independent of $D_1,D_2,\epsilon$ introduced below.
All remaining contributions on $\mathcal F_x^c$ are bounded by
actual--clean child discrepancy norms with an amplification at most
$2^{C d}$ at slack $d$.
\end{lemma}

\begin{proof}
For each small dyadic group, let $\Delta_d^{\mathrm{cl},\nu}(x,y)$ be
its complete clean true-ID increment in channel $\nu\in\{1,2\}$,
including any exact base term assigned to that group.  Take an absent second channel
to be zero, and choose a deterministic descriptor bound $k_d\le2^{Cd}$.
By Lemma~\ref{lem:est-raw-difference}, the clean paired increments at a
group of slack $d$ obey their crude absolute estimates in the retained
endpoint metric.  These estimates hold for each fixed valid variant
after its measure-preserving change of variables.  Counting variants
and summing difference layers by the triangle inequality costs
$2^{O(d)}$, without using walk contraction.  Within each paired layer
the true-ID cancellation is performed before this inequality is
applied.  If a local unpaired base term occurs in a small group, its
first-group slack satisfies $M=O(d)$, so the same crude bound holds in
its normalized raw units.  Since $\sup_y|Z_y|^s\le\sum_y|Z_y|^s$ over actual ID
positions, the row norm of
$\max_\nu\sup_y|\Delta_d^{\mathrm{cl},\nu}(x,y)|$ is at most
$2^{-p+O(d)}$.

Lemmas~\ref{lem:buf-update} and \ref{lem:buf-clamping}, together with
Equation~\eqref{eq:buf-propagation}, propagate an input discrepancy at
slack $d$ to the final array with a factor polynomial in $d+1$.
At each step the discrepancy is measured between the clamped full
arrays, including their defaults.  The next bucket increment is
compared by Lemma~\ref{lem:fp-buckets}; compression then applies its
uniform Lipschitz bound and the common clamp cannot enlarge the
discrepancy.  This recurrence does not require their later bucket
distributions or stored lists to agree.
Root-Lipschitz row functions and bounded final support then give the
same type of bound for the output metric.  Choose fixed $C_P,A_P$ so
that $C_P(d+1)^{A_P}$ bounds this propagation.  For a retained row define
\[
 H_x=C_P\sum_{\substack{d\le2h_*\\d\text{ a group slack}}}
       (d+1)^{A_P}(k_d+1)
       \max_\nu\sup_y|\Delta_d^{\mathrm{cl},\nu}(x,y)|.
\]
Set $H_x=0$ when $X_u(x)=0$.  This quantity bounds the propagated clean-height terms in
Equation~\eqref{eq:fp-bucket-bound}.  The preceding estimates give
\[
                 \|H\|_{\mathrm{row},u;s}\le2^{-p+C_0h_*}.
\]
All increments defining $H_x$ are clean, and its other factors are
deterministic.  Thus $H_x$ and $\mathcal E_x$ are
$\mathcal G_b$-measurable.  Conditioning on this sigma-field gives
\[
 \mathbb E\bigl[\mathbf1_{\mathcal C_x}H_x^s\mid\mathcal G_b\bigr]
 =H_x^s\Pr(\mathcal C_x\mid\mathcal G_b)
 \le C\,2^{2C_Eh_*}B^{-1/2}H_x^s.
\]
The row membership $X_u(x)$ is also $\mathcal G_b$-measurable.
Integrating the last inequality in the row measure therefore yields
\begin{align*}
 \|\mathbf1_{\mathcal C}H\|_{\mathrm{row},u;s}^s
 &\le C\,2^{2C_Eh_*}B^{-1/2}
              \|H\|_{\mathrm{row},u;s}^s.
\end{align*}
Taking roots proves Equation~\eqref{eq:fp-direct-charge} after enlarging
$C$.  We have not conditioned on $\mathcal F_x^c$ in this independence
calculation: the actual error there is bounded by
$\mathbf1_{\mathcal C_x}H_x$, which is then integrated on the whole
space.  Conditioning on the salt-dependent stay-inside event is
unnecessary and could destroy the stated independence.

For the first term of Equation~\eqref{eq:fp-bucket-bound}, split products
using bounded-factor differences, actual or clean row caps, and the
actual or clean incoming cap of the represented two-way factor, exactly
as in Lemma~\ref{lem:est-raw-difference}.  These are comparisons of
unaliased sums using the respective children; they make no assumption
about key success or independence of child errors.  Fixed valid
variants, scale costs, and counts are all at most $2^{O(d)}$, with the
possible source-release factor absorbed by the same precision gain as
in Equation~\eqref{eq:fp-actual-incoming}.  Compression propagation is
again absorbed in $2^{C d}$.  This proves the remaining assertion.
\end{proof}

\Needspace{8\baselineskip}
\subsection{Close the accuracy and comparison bounds together}

There are three assertions to prove simultaneously: accuracy of the actual estimator, accuracy of the exact-equality comparison, and closeness of the two. Lower-block calls coincide in the comparison. Within the current block, the induction permits a comparison error that grows with the budget, leaving a geometric gain whenever a child has lower budget. That gain absorbs the cost of enumerating paths and propagating local errors.

\begin{proposition}[Uniform block coupling]\label{prop:fp-coupling}
The fixed constants can be chosen so that, for every task of budget
$M\in\mathcal I_b$ and precision $p$, each prescribed actual--clean
output discrepancy, including every pre-incoming score and witness,
is at most
\begin{equation}
 2^{-p}\lambda_M,\qquad
 \lambda_M=2^{-D_1h_*+D_2(M-M_b)}.
 \label{eq:fp-lambda}
\end{equation}
The same assertion holds for the ordinary vector metrics.  Adding
these discrepancies to the clean estimates proves the target error
bounds for the actual tables and vectors.
\end{proposition}

\begin{proof}
Perform strong induction on the statistical budget, simultaneously for
the actual target estimates, the block-clean target estimates needed at
that budget, and Equation~\eqref{eq:fp-lambda}.  An inactive task returns
the same zero outputs and empty paths in both versions, so its coupling
error is zero.  Lower-block children are identical in a block coupling
and already have their actual target bounds.  Same-block clean children
have the clean target bounds at smaller budgets.  Therefore
Lemma~\ref{lem:fp-clean}, the prefix estimate, and the compression
analysis give the clean active target bounds with the spare constants
in Proposition~\ref{prop:est-clean}.

For the coupling, Lemmas~\ref{lem:fp-escape}--\ref{lem:fp-direct}
bound the pre-incoming quantities.  Incoming gates are then compared
by their rank and root-Lipschitz bounds.  In their multiplier effect,
restrict to the union of the final supports of the two tables; these
have bounded incoming multiplicity.  Thus Lemma~\ref{lem:atom-assembly}
costs only a fixed factor on pre-incoming errors.  Its separately queried
column vectors have a positive constant budget drop and their
$f/u$ metric factor is covered by their precision gain.  Vector raw
sums have no endpoint grouping, but their same-block children may
differ.  The same crude individual-variant comparisons and nonexpansive
clamps bound those differences.  Detector rank comparisons are already
part of the assembly estimate.  No simultaneous key-success event over
an incoming search is required.

Collect the finite number of output types and absorb their fixed
assembly factors into one fixed exponent $C$.  For a child budget
$m\in\mathcal I_b$ used at group slack $d$, we have
\begin{equation}
 \lambda_m=\lambda_M2^{-D_2(M-m)}
              \le\lambda_M2^{-D_2d}.
 \label{eq:fp-child-gain}
\end{equation}
For a separate post-assembly query with constant positive budget drop,
use $d=1$ in this inequality.  Replacing the dyadic group sum by the
larger sum over positive integers, all coupled child charges are at most
\begin{equation}
 2^{-p}\lambda_M\sum_{d=1}^{2h_*}2^{-(D_2-C)d}.
 \label{eq:fp-recurrence}
\end{equation}
Counts of budget layers and dependency types are included in $2^{C d}$.
Groups with larger slack have identical children and contribute zero.
Together with Equation~\eqref{eq:fp-direct-charge}, this is the complete
coupling recurrence.

We can now choose constants in the same order as the two error sources. First make lower-call discrepancies contract across a budget decrease; then make direct collisions rarer than the smallest allowed comparison error in a block. Choose $D_2>C$ so large that the infinite geometric sum in
Equation~\eqref{eq:fp-recurrence} is at most $1/4$.  Then choose
$D_1>D_2+1$.  Finally choose rational $\epsilon>0$ so small that
\begin{equation}
                  \frac1{2s\epsilon}>C+D_1+2.
 \label{eq:fp-epsilon}
\end{equation}
Since $h_*\le\epsilon\log_2B$, the direct relative charge satisfies
\[
 2^{C h_*}B^{-1/(2s)}
 \le 2^{-(1/(2s\epsilon)-C)h_*}
 \le\tfrac14\,2^{-D_1h_*}
 \le\tfrac14\lambda_M,
\]
after increasing the fixed minimum $B$ to absorb any fixed coefficient.
The child and direct charges fit within Equation~\eqref{eq:fp-lambda},
even with room for fixed norm conventions already included in $C$.

Moreover $0\le M-M_b<h_*$ implies
\[
              \lambda_M\le2^{-(D_1-D_2)h_*}.
\]
This is uniformly as small as required after the same minimum-size
padding.  The triangle inequality therefore adds the coupling bound to
the spare clean active bound, for example $C_b2^{-p}/2$, without
exceeding $C_b2^{-p}$.  Inactive target bounds are the previously fixed
bounded-range bounds.  This completes all three inductions.

The order of choices is consequential.  Fix the analytical hierarchy,
table thresholds, and trivial-case bound first.  Choose the precision
and cutoff constants for the mean comparisons; then the compression
floors, capacities, and walk lengths so the clean sums converge with
their spare coefficients.  Fix clamping ranges and the crude exponents
in the fixed-variant comparisons next.  These choices determine $C$.
Only after them choose $D_2,D_1,\epsilon$ as above.  Thus reducing
$\epsilon$ to obtain the direct collision gain does not alter an
earlier clean-coefficient requirement.

Later choices of the common stage-size bound, top-budget coefficient,
and environment-field exponent do not change this crude exponent $C$.
Immediate descriptor counts depend on relative slack, fixed support
bounds and fixed walk length.  The sampler gap is uniform in both
field size and the length of its single free-bit component.  Likewise,
an enlargement of the fixed identifier-length coefficient only changes
the polynomial-evaluation collision term in Lemma~\ref{lem:fp-hash} to at most
$2c_{\mathrm{ID}}/B^2$; increasing the fixed minimum
$B\ge c_{\mathrm{ID}}$ bounds it by $2/B$.  These enlargements can
therefore be absorbed after $D_2,D_1,\epsilon$ have been fixed.
\end{proof}

\begin{theorem}[Error bounds for the actual estimators]
\label{thm:fp-estimators}
For every allowed task, with precision $p$ defined in
Equation~\eqref{eq:est-precision}, the actual row outputs, including the
pre-incoming scores and support witnesses, satisfy
\[
                \mathfrak d_{u,a}(\bar H,H)\le C_b2^{-p}.
\]
The clipped vectors have the corresponding ordinary error bound.
In particular, every allowed reward query satisfies
\begin{equation}
 \frac1N\mathbb E_\sigma\sum_{x\in V_l}\frac{X_u(x)}u
           |\bar W_l(x)-W_l(x)|^s\le C_b^s2^{-sp}.
 \label{eq:fp-reward-bound}
\end{equation}
These estimates are uniform over the allowed task parameters and stage,
and the constants are independent of the input and of $B$.
\end{theorem}

\begin{proof}
Proposition~\ref{prop:fp-coupling} adds the actual--clean discrepancy
to Proposition~\ref{prop:est-clean} with its spare constant.  The
inactive-task bounds follow from the fixed bounded ranges and supports.
The induction in Proposition~\ref{prop:fp-coupling} includes both kinds
of tasks and every prescribed output.  Applying its ordinary vector
statement to $W_l$ and taking the $s$-th power gives
Equation~\eqref{eq:fp-reward-bound}.
\end{proof}

